% Completed guide assembled as the successive parts of the cote24 audit were finished.
\subsection*{Piotr Blass's translation: extent, opening, and notation}
The 1987 typescript consists of a title page, Summary, contents, and the English text of IV 20, Sections 1--11, on archive sites 4--74. Its contents list further sections, but the translated mathematical body ends with the Grassmannian discussion corresponding to French site 65. The handwritten French note on site 75 is retained in all three Blass readers; the stray CV sheet on site 76 is retained only in the diplomatic reader. These books preserve Blass's English, his French-language glosses, his typed questions, and the visible handwritten annotations. The following comparisons come from the two checked witnesses, and the decision records locate each gap or query by its French archive site. Sources: \href{https://grothendieck.umontpellier.fr/web/viewer.html?file=../24.pdf\#page=2}{cote 24, sites 1--76}; \href{https://grothendieck.umontpellier.fr/web/viewer.html?file=../25.pdf\#page=3}{cote 25, IV 20}.

The English title page identifies the work as ``Stage 0 -- translation by Piotr Blass,'' and the following Summary places the translated \emph{Commentaires} before the contents. Its queries belong to this witness: ``reorganize???'' corresponds to \emph{réorganisation de fond en comble} on French site 2; the typed question about no.\,16 and its handwritten ``Ask AG.'' refer to the French statements on sites 2--3. In Section 1, the English develops the same quotient-Module, hyperplane, relative-divisor, and incidence constructions as French sites 4--8. The opaque strips on English sites 4 and 6 remain unreadable, with their positions identified rather than completed from the French. English site 5 asks which paragraph 17 is meant and queries the separate circle above the functor domain; both corresponding passages are on French site 5. English site 7 retains four unfilled arrow positions and asks whether \emph{Module} always means a coherent or quasi-coherent sheaf of modules, at the passage corresponding to French site 7. Its drawn diagram contains four connectors without arrowheads. The comparison also records exact differences of notation: the final term of the transpose formula on English site 4 is $\underline O_P(1)$, while French site 4 types $\underline D_P(1)$ and then names a section of $\underline L_P(1)$; English site 4 types an ordinary $p*$ where the French has a handwritten subscript star; English site 5 lacks the French dual marks in the symmetric-power expression and on the final projective space; English site 6 lacks one dual mark in the section over $S'$; and English site 7 types the continuous string $XPV$ where French site 7 says that $F$ is a Module on $X$. These visible forms are preserved. Full-height baseline letters and genuinely raised characters have been compared separately, since the typescript uses both. Sources: \href{https://grothendieck.umontpellier.fr/web/viewer.html?file=../24.pdf\#page=2}{cote 24, sites 1--8}; the corresponding French passages are \href{https://grothendieck.umontpellier.fr/web/viewer.html?file=../25.pdf\#page=3}{cote 25, sites 2--8}.

\paragraph{Blass's Section 2 (sites 9--19 opening).}
The translation follows the local study from irreducibility and finiteness through associated cycles, regularity, smoothness, and separability, and retains the author's comments on reorganizing the results. Its queries can be located precisely: site 9 questions the insertion \emph{En tous cas} in French 8; site 10 asks about underlined capitals; site 12 queries the added alternative in Proposition 2.6 and the paragraph-3 reference of French 9--10; site 13 leaves the marginal remark about the \emph{dim.\,|\,dém.} of 2.1 unread and retains four reference placeholders from French 11; site 15 leaves the argument of $k(\ )$ blank, as French 13 does; and site 17 asks about the two paragraph-6 references on French 13--14. Both opaque masks on English 16 remain unreadable. The first ``PB incomprehensible'' on English 18 corresponds to the repeated $Z_i$ in French 14's first assertion of 2.16; the repetition is present in both witnesses. Its second query accompanies the reduction to an algebraically closed field, whose order relative to the separability sentence differs from the French insertion. The mathematical comparisons retain all observed differences: English 9 inserts $k$ in $\operatorname{Spec}k(\eta)$, where French 8 has $\operatorname{Spec}(\eta)$; English 10 types $Z_\eta$ in the proof where French 8 corrects the letter to $Y$; English 11 types $Z1$ where French 9 has $Z'$; English 13 types a capital $N$ in the coprofundity bound and $H_\eta$ in the conclusion, where French 11 has lowercase $n$ and the doubtful $E|F$ with indices $\eta,y$. English 14 places its first bar over \emph{dim}, omits the inverse-image exponent in $h(Z)$, and says that constructibility is already known by applying 2.3, while French 11--12 bars $f(Z)$, writes the inverse-image exponent, and says that constructibility must be known to apply 2.3. English 15 has a lowercase $y$ and ordinary arrow into $U_{k(\eta)}$ against French 12's capital $Y$ and hooked arrow, and retains the baseline $1$ and lowercase $u/k$ in its differential notation. English 16 places the bound $1\leq i\leq n$ beneath the quotient sum; French 13 places it after the inserted generators $a_i$. English 16--17 omit the French marginal qualification beside 2.14, whose first term remains doubtful, while English 17 types $k(y)$ in the equivalence and omits ``geometrically'' before the first ``normal'' in 2.15. Capital $K(\eta)$ occurs on English 16--17 where the corresponding French uses lowercase $k$. Several additional underlines on $Y$ are visible on English 9, 11, 12, and 14. English 18 replaces the French ranges 2.13--2.17 and 2.12--2.17 by conjunctions of their endpoints; the explanation continues on English 19, with an extra typed arrow before ``smooth = regular.'' None of these forms is repaired from the other witness. Sources: \href{https://grothendieck.umontpellier.fr/web/viewer.html?file=../24.pdf\#page=10}{cote 24, sites 9--19}; \href{https://grothendieck.umontpellier.fr/web/viewer.html?file=../25.pdf\#page=9}{cote 25, sites 8--15}.

\paragraph{Blass's Section 3 (sites 19--21).}
The translation presents the Bertini--Zariski theorem through the function-field description and Zariski's lemma, retains the request that Dieudonné choose a proof from the literature, and follows with the geometric-integrality and connectedness corollaries. The numbering query on site 19 corresponds to the overwritten heading on French 15; ``[PB: check this!]'' points to the inclusion described by sending $S_i$ to $S_i$ on French 16. Several literal formulas require care. English 19 has $\dim\overline{f(X)}\geq0$ where French 16 has $>0$, and ends $L=K(S_1,\ldots,S_n)$ at $n$ where the French ends at $n-1$. English 20 uses $a_iS_1$ where French 16 has $a_1S_1$, and omits the minus before $t_it_n^{-1}$ in the definition of $a_i$. Its lemma omits the French definition $L=K(S_1,\ldots,S_m)$, while the sum has printed capital limits $I$ and $M$ against French 17's $1$ and $m$; the English field name is the literal baseline $k1$. English 21 uses a clear typed $\sigma$ and $f(X)$ where the checked French 17 retains doubts about the script family letter and the argument $Z|X$. It also prints capital $K(\eta)$ in the conclusion and ``$X_0$ and $X'$'' at the start of the chain, against the French lowercase field letter and equality. Its handwritten query points to the paragraph-6 reference accompanying Hartshorne's theorem. The guide reports these readings and retains the French uncertainties. Sources: \href{https://grothendieck.umontpellier.fr/web/viewer.html?file=../24.pdf\#page=20}{cote 24, sites 19--21}; \href{https://grothendieck.umontpellier.fr/web/viewer.html?file=../25.pdf\#page=16}{cote 25, sites 15--18}.

\paragraph{Blass's Section 4 (sites 21--25).}
The translation returns to the general base and uses constructible properties to pass from the generic hyperplane to sufficiently general hyperplanes, then lists the Module and prescheme properties of 4.2--4.4. The doubtful handwritten query on English 22 points to the paragraph-9 reference in French 18; English 25 asks both what AQT means and where paragraph 9 is, at the opening of French 21. English 22 prints $G$ at a hypothesis where French 18 retains the letter doubt $F|E$, omits ``two'' from ``the following sections,'' and inserts a further $\xi$ before ``$k$ algebraically closed.'' English 23 uses ``scheme,'' \emph{coprof}, and ``Cohen Macauley'' where French 19 has \emph{préschéma}, \emph{coppof}, and ``Cohen-Macaulay''; its characteristic character is the round letter O, distinguished on the same page from the taller numeral 0. The wording ``F is satisfied at P'' is also literal. English 24 lists geometric integrality before geometric irreducibility, following the direction of the French relocation while the overwritten French labels remain uncertain. More consequentially, its characteristic-zero and unramified clauses are joined by ``and,'' against French 20's \emph{ou}; the French inserted morphism description itself still carries a reading doubt. English 25 prints $K(\xi)$ against lowercase $k(\xi)$, and $f(X)$ against $f(X_i)$. Its $E_\xi$ and final projective-space symbol without a dual are already present in French 20. The comment following 4.4 says ``not valid as such if f is quasi-finite,'' whereas French 21 says \emph{n'est valable tel quel que si f est quasi-fini}; the next instruction becomes ``it is sufficient to consider'' where the French says \emph{il faut regarder}. The comparison records these exact changes in qualification without repairing the statements. The Latin instruction is retained with Blass's added label ``[Latin].'' Sources: \href{https://grothendieck.umontpellier.fr/web/viewer.html?file=../24.pdf\#page=22}{cote 24, sites 21--25}; \href{https://grothendieck.umontpellier.fr/web/viewer.html?file=../25.pdf\#page=19}{cote 25, sites 18--21}.

\paragraph{Blass's Section 5 (sites 25--29).}
The opening incorporates the paragraph on the openness of $U$ and $E$ that French 21 cancels and French 22 explicitly restores as text struck out in error. The query ``PB ???'' and the handwritten ``illegible'' on English 25 refer to its definition of $\xi$ as the image of $y$; English 26 then completes the paragraph and develops the section defining the universal hyperplane divisor. The following regularity and openness statements correspond to French 22--24. English 26 leaves a star after ``Mike'' without a footnote and question punctuation in the parenthesis about $S'$ over $S$; French 22 has \emph{comme dirait Mike} and canceled doubtful question marks in that parenthesis. English 27 leaves a question at $\mathrm{id}_F\otimes$ where French 22 has $\mathrm{id}_F\otimes\varphi$, and a further star at the paragraph-3 reference of French 23. Its incomplete $o_{\mathrm{III}}$ reference corresponds to French $E_{\mathrm{III}}$. In the proof of 5.3(b), English 27 prints equality between the intersection and the repeated $Z$ term, against French 23's handwritten inequality; its final section symbol also lacks the French $\eta$ index. English 28 has ``P is flat over S'' against the inserted French ``F,'' and ``a little less evidently'' against \emph{évidemment}. It asks for the references 12.1.1 and 12.1.4 cited at the start of French 24's 5.6. English 29 leaves blank the relation between $U_P$ and $U$, where French 24 inserts inclusion, omits $\xi$ after the parenthetical $Y$, and lacks $\varphi_\xi$ in the regularity condition of 5.7. The incomplete bracket ``[fascicule'' also remains. Sources: \href{https://grothendieck.umontpellier.fr/web/viewer.html?file=../24.pdf\#page=26}{cote 24, sites 25--29}; \href{https://grothendieck.umontpellier.fr/web/viewer.html?file=../25.pdf\#page=22}{cote 25, sites 21--24}.

\paragraph{Blass's Section 6 (sites 30--31 opening).}
The translation combines the generic connectedness statement 3.3 with the cited 15.6.3 in Proposition 6.1 on arbitrary hyperplane sections, and retains the planning remarks about 4.4, sections of excessive dimension, and grouping the connectedness questions in one number. The query ``[x X proper over k suppose]*'' on English 30 corresponds to the insertion ``X propre sur S'' in French 25; there is no starred footnote on that English page. English 30 leaves $f(Z)$ unbarred where French 26 has a handwritten closure bar, and prints ``if Y is proper over k'' where the French gives \emph{car Y est propre sur k}. Its three long diagram arrows are drawn in pen and are marked accordingly in the checked transcription. English 31 adds ``(or section)'' after ``No.'' in the planning sentence. Sources: \href{https://grothendieck.umontpellier.fr/web/viewer.html?file=../24.pdf\#page=31}{cote 24, sites 30--31}; \href{https://grothendieck.umontpellier.fr/web/viewer.html?file=../25.pdf\#page=26}{cote 25, sites 25--26}.

\paragraph{Blass's Section 7 (sites 31--35).}
This section follows the use of rational points over infinite fields and étale base changes over finite residue fields to construct hyperplane sections and multisections, then the statements 7.2--7.4 and Remark 7.5. It retains the references to 4.1--4.3 and no.\,5 and the induction on fiber dimension in 7.4. English 31 types the literal section sign \S\ for the section variable where French 27 has $\xi$, and initially has $\xi$ without the French inserted index $0$, although $\xi_0$ appears later in the English paragraph. English 32 asks about ``closer and closer (by successive approximations ???)'' for French 28's \emph{de proche en proche}, with a handwritten ``Translator's note.'' Its printed $k1$ lacks the denominator of the French residual extension $k'/k$. English 33 repeats ``finite'' in the presentation-and-flatness hypothesis and uses $\xi,s$ as the indices of $Y$ and $F$, against the French $\xi,\xi(s)$. The neighborhood in $E$ occurs in both witnesses; the French range in the reference to 4.3 remains doubtful. English 34 leaves several placeholders at French 29's added irreducible-component condition and the references 12.2.1(x) and (xi); part of the French added N.B. also remains uncertain. ``Text crossed out'' represents the canceled paragraph continuing into French 30's Stein-factorization discussion. The underlined ``X [blank] U'' query in 7.4 remains unresolved by French 30, which likewise shows no intervening sign. The raised $1$ in a later $S^1$ is retained alongside baseline $S1$ occurrences. English 35 omits ``over'' before the final $S$ in the nonprojective case and leaves ``ILLEGIBLE'' for the French marginal suggestion about avoiding an \emph{ensemble} whose adjective remains doubtful between \emph{fini} and \emph{donné}. Sources: \href{https://grothendieck.umontpellier.fr/web/viewer.html?file=../24.pdf\#page=32}{cote 24, sites 31--35}; \href{https://grothendieck.umontpellier.fr/web/viewer.html?file=../25.pdf\#page=28}{cote 25, sites 27--30}.

\paragraph{Blass's Section 8: exceptional hyperplanes (sites 36--42).}
The section begins with the dimensions of exceptional hyperplane loci, then the coprofundity statements and the regularity criteria 8.4--8.6. English 36 repeats the clause defining the set of $\xi$ and says ``a better upper bound'' where French 31 changes \emph{majoration} to \emph{minoration}; the query asking for a number corresponds to the French's own unfinished \emph{cf n°}. English 37 asks whether a numerical hypothesis is two, at French 32's $\dim f(X_i)>0$, and uses $G_{k(\xi)}$ in its first coprofundity expression where French 32 has $F_{k(\xi)}$. It omits the intervening assertion about $\operatorname{coprof}G_\xi$ and prints an incomplete $O_{\mathrm{IV}}$ reference against French $E_{\mathrm{IV}}$. Its later gap corresponds to French 33's assertion about $G_\xi$ satisfying $(S_{k-1})$. The extensive queries on English 38 point to French 33's partly unreadable opening, its clear $k=2$, the added reference to 8.5, and the unfinished handwritten discussion at the foot. English 38 says the $(S_2)$ conditions are met where the French says \emph{ne sont pas}. English 39's question about an insertion follows the component clause developed in French 34's relocated continuation, with its zero-dimensional image and quasi-finite case; its lemma retains the literal bound $\operatorname{codim}(Y,X)\geq K$. English 40 leaves blank the relation that French 35 gives as $\operatorname{codim}(T_i,X_j)=k$ and omits the intervening alternative involving $\dim f(X_j)>0$, a point outside $Z_0$, and the next exceptional clause. Its local-ring indices $Xj,Ti$ share one lowered tier, and the production layer of the inline slashed C remains uncertain. English 41's gaps and ``floppy or perhaps sloppy'' comment accompany an already incomplete French 35 statement; its later ``R of'' gap corresponds to French 36's $X$, and its ``finite'' differs from French \emph{quasi-fini}. English 42 omits the dimension parenthesis, \emph{est vide}, and a displayed coprofundity bound from French 36, then prints ``by preference'' for \emph{à priori}, $Z=X$ for $z=x$, and a depth bound $\geq k$ where the French has $<k$. Its final editorial sentence changes the instruction about stating a corollary with 8.6, and the new differential discussion begins with $f:X\to f$ against French 37's $f:X\to P$. The footer query has no matching body star, so its exact scope remains unsettled. Sources: \href{https://grothendieck.umontpellier.fr/web/viewer.html?file=../24.pdf\#page=37}{cote 24, sites 36--42}; \href{https://grothendieck.umontpellier.fr/web/viewer.html?file=../25.pdf\#page=32}{cote 25, sites 31--37}.

\paragraph{Blass's Section 8: conormal calculations and 8.7--8.9 (sites 43--50).}
The conormal Module, universal singular-zero locus, and incidence interpretation lead into the differential diagram and Theorem 8.7. The query about the last two symbols on English 43 points to the kernel of $E_X\to O_X(1)$ on French 37; the queries on English 44 point to French 38's clipped \emph{comm}, \emph{ledit sous-espace vectoriel}, and linear subvarieties containing $a$. The conormal letter retains its unresolved alphabet, and the English section symbol after ``image by'' lacks the French prime. English 45 prints $dt(r-d-1)$ where French 39 has $d$ followed by a doubtful plus, asks about paragraph 16, and leaves a query after the differential-section formula $\psi=d\phi|Y$. Both English $O$ indices are lowercase $y$, the second underlined; both $P$ letters in the first $\Omega$ index occupy a deeper lowered tier. English 46's diagram has typed objects and pen connectors, including one connector without a secure head and two dashed diagonals. Its $B$ and partial-shaped labels differ from French 40's $\beta$ and $\alpha$; its top object, the copied uncertain pen letter, and the missing lower-left base letter remain unresolved. The question about $\psi$ corresponds to French 40's inserted $\mu$. Its displayed composition goes directly to the $(1,1)$ target, whereas the French first writes $(1,0)$ and then interprets the composition as a section with $(1,1)$. English 47 retains the duplicate ``above lies over'' and joins the announced result to the next editorial paragraph. English 48 preserves the paragraph-16 and definition-reference queries and leaves the first projection unread, where French 41 gives $x$; its $H$ bears a raised numeral 1 against the French prime. More substantially, English 48 jumps from ``image of the tangent map of'' to ``the closure of its image,'' omitting the end of 8.7(c), its displayed map, the proof and diagram (D'), and the heading and opening hypotheses of Corollary 8.8 on French 42. The later conjunction ``(i) and (ii) (bis)'' replaces the French range. English 49's queries correspond to French 43's N.B. about the immersion hypothesis and the inserted \emph{ou}; it prints $\leq r-2$ where the French changes the sign to equality, $T$ for $Y$ in the image clause, twice ``(ii ter)'' for ``(ii bis),'' unbarred capital $K$ for barred $k$, and ``(i leir)'' for ``(i bis).'' Its later ``(i)'' corresponds to a doubtful French (ii)/(iii), and the proof lacks the French point $y$. English 50 repeats the properness and surjectivity clauses, omits the French qualification about finding an open subset, and says ``to T and $\xi$'' for \emph{en $\xi$}. Its $H1$, four $g-1$ groups, and two $Y1$ groups remain on the baseline, alongside genuine $g_1$ indices. The query after $h$ corresponds to canceled fragments on French 44; English $\psi$ differs from the eventually secure French $\mu$. English 50 bars the second $Y1$ in the closure sentence, against the first $Y^1$ in French 45. Sources: \href{https://grothendieck.umontpellier.fr/web/viewer.html?file=../24.pdf\#page=44}{cote 24, sites 43--50}; \href{https://grothendieck.umontpellier.fr/web/viewer.html?file=../25.pdf\#page=38}{cote 25, sites 37--45}.

\paragraph{Blass's Section 8: singularities and linear subsystems (sites 51--62).}
The continuation completes 8.9, discusses the singular and supersingular loci and 8.10--8.18, then returns to linear subsystems and the projected construction. English 51 prints baseline $U1$ for French $U'$, omits the explicit isomorphism induced by $g_1$, and has equality in the dimension bound where French 45 retains a doubtful replacement by $\leq$. Its query about the singular locus points to the French definition by the vanishing of a section $D$ of an invertible Module. English 52's queries refer to French 46's tangent at $x$, its revised singular-locus notation, and the repeated geometric-singularity phrase; the English adds ``double'' to ``ordinary points,'' and its clear degree 4 leaves the French 4/5 doubt unresolved. English 53 retains its literal $W$ formula and index levels against the French handwritten omega. Its three separate gaps point to French 47's top note and the inserted dimension-one example; the French marginal definition of \emph{réglé} is absent. English 54's hyperplane query corresponds to French 47's reorganization note, and it omits the live premise that $X$ satisfies geometric $(R_k)$. In 8.13 it prints capital $K$ and the bound $\leq K-1$, against French 48's $\leq k+1$, while leaving out the marginal geometric-singular-locus and dimension notes. English 55 retains a repeated depth expression, $\operatorname{prof}xX\geq2\operatorname{prof}_xX$, against the single French bound, and omits ``geometrically'' before ``normal'' in 8.15. English 56 leaves the codimension bound blank where French 49 has $\geq2$, retains a narrow left angle after $Z$, and prints $K+1$ against a doubtful French letter or digit. English 57's four gaps correspond individually to French 50--51; 8.17(a) omits the French qualification beginning \emph{Il est possible que} and the geometric-irreducibility phrase, prints $f(x)$, and calls the tangent hyperplanes ``hyperplane sections.'' English 58 preserves the French-language query about \emph{réglé}, ``non???'' before ``singular,'' and a further gap of uncertain attachment. It omits $d$ before the dimension bound, uses a slash for restriction, and raises 1 where French 52 has a prime. English 59 omits two image-closure bars and dual marks, prints capital $X\in T$, and says ``and not only'' where the French says \emph{mais seulement}; its translation alternatives and ``Fr p50'' stay literal. English 60 preserves further alternatives and French glosses, omits several dual, zero, and $\xi$ marks, prints $f(x)\ C\ L0$ against French 53's $f(X)\not\subset L^o$, and uses capital $U$ and raised 1 in the diagrams. English 61 prints $\dim f(X)=2$ where French 53 has a not-equal sign and a doubtful 2/1, omits field bars and the zero on $L$, starts the field expression with capital $K$, says ``dimension at most one'' for \emph{dimension une de plus}, and leaves the target image unbarred. The final English 62 remark retains lowercase $x$ where French 54 has $X$, and its ``Bertinisque'' spelling. Every remaining query is kept with its recorded French location. Sources: \href{https://grothendieck.umontpellier.fr/web/viewer.html?file=../24.pdf\#page=52}{cote 24, sites 51--62}; \href{https://grothendieck.umontpellier.fr/web/viewer.html?file=../25.pdf\#page=46}{cote 25, sites 45--54}.

\paragraph{Blass's Section 9 (sites 63--66).}
The translation follows the Veronese-type constructions and the degree-dependent results through the proof of 9.3, then preserves the unfinished 9.4 and 9.5. English 63 retains the missing relation in the opening direct-image/Sym expression, also absent in French 55; its ordinary and lowered stars and its question marks are separately preserved. Its literal $O\ P(n)$ uses ordinary baseline letters, and the following $P(n)$ lacks the French dual mark. The wide gap after ``embarrassing f'', the French phrase without accents, and ``or (admits)'' correspond to French 55's \emph{de f à certains égards} and \emph{ou n'admet qu'un}. English 64 repeats an N.B. clause that appears once in French 55--56. In 9.3(c), ``where'' corresponds to the inserted French \emph{ou}; the field is unbarred, the first $Y(n)$ lacks the attached index, and the last exponent has a closing parenthesis without its opening mate. English 65's queries correspond to French 57's clause about $t_s$, \emph{on fait pareil}, and a live/canceled-text boundary. The affine-model sentence uses $f(X)$ and ends its coordinates at $T_s$, against French 56's $f(x)$ and a doubtful $T|X$ with index $r$. ``Move margin'', capital $K$, and the capital C in the image expression are literal. The $i/t_s$ mismatch is shared with French 57. English 66 prints a clear $Q/R$ where the French retains $Q/P$ doubt. Its blank Corollary 9.4 and all the 9.5 questions point to French 57's unfinished headings, conditions, and notes. Neither witness supplies a completed statement there. Sources: \href{https://grothendieck.umontpellier.fr/web/viewer.html?file=../24.pdf\#page=64}{cote 24, sites 63--66}; \href{https://grothendieck.umontpellier.fr/web/viewer.html?file=../25.pdf\#page=56}{cote 25, sites 55--57}.

\paragraph{Blass's Section 10 (sites 66--71).}
The section translates the discussion of sections by linear subvarieties, the regular-immersion construction, and the planning comments around 10.2--10.4. English 67 keeps ``not meeting k'' and adds ``(Z AG's error P.B.)'' with a handwritten question; French 58 indeed has $k$, so this remains Blass's proposed correction. Several closure bars and dual marks are absent, the English uses raised 1 for a French prime, and its clear $f(T)$ leaves the French $F|T$ unresolved. English 68 keeps ``pull out of find out'' and two N.B. closures. Three sheaf occurrences have the narrow digit $0$ where other occurrences on that page use round capital $O$; the underlines, product letter x, and typed dual v's remain as printed. English 69 refers to ``No.1'' where French 60 inserts 12, repeats overlapping opening clauses, and retains the mystery comment, the gap before 5.3, and the handwritten request. The corresponding French reference remains doubtful between S.3 and 5.3. Its ``only'' replaces \emph{surtout}, ``which according to 10.2'' replaces \emph{qu'après 10.2}, and the codimension-two wording lacks the uncertain sign visible in French 60. English 70 preserves the page reference ``[58]'', the extra qualification ``to some extent'', and ``except'' in a sentence whose French has repeated \emph{sont}. The query about a regular-immersion section points to French 61's \emph{un n° en rabiot}. The typescript then says ``Proposition10.3 is crossed out. AskAG if that is his intention.'' It omits the proposition body, the exercise paragraph, and the dimension-at-most-one transition present in checked French 61; the claim about cancellation is attributed to Blass. English 71 asks for the translation of ``sorites'' and adds ``(facts)'', at French 62's doubtful \emph{sorties|sorttes}; its handwritten translation question remains partly uncertain. Other marked alternatives concern taking the construction in practice, revisiting it, $X$ or $S$, and closed immersions. ``Forefathers'' renders \emph{Papa}, the phrase \emph{au dessus de s} is omitted, and $n-1$ corresponds to a doubtful French $t|r-1$. Sources: \href{https://grothendieck.umontpellier.fr/web/viewer.html?file=../24.pdf\#page=67}{cote 24, sites 66--71}; \href{https://grothendieck.umontpellier.fr/web/viewer.html?file=../25.pdf\#page=59}{cote 25, sites 58--62}.

\paragraph{Blass's Section 11 (sites 71--74).}
The final translated section develops Grassmannians and linear subspaces of projective bundles, corresponding to French 62--65. English 71 prints $nE$, a plain rank-one $G$, and $P(nE)$ without the French exterior-power signs, and leaves several sheaf letters ununderlined. English 72 places handwritten $r$ above typed $(n)$ with only $n$ inside the parentheses; it repeats two representability lines where French 63 has one clause. It twice prints capital $S$ in the Grassmannian index, says ``recovered from'' for \emph{recouvert par}, and omits ``vector'' from ``fiber bundle''. The sequence endpoints are round capital O rather than the French numerical zeros; all four English connectors have arrowheads, where the two middle French strokes do not. The Hom subscript is a narrow digit 0 followed by lowercase s on one lowered tier, against French underlined O with lowered S. English 73 omits the characterization by closed subpreschemes that are smooth over S and everywhere of the stated relative dimension, joining directly to projective degree one. Its ``empty is n=0'' corresponds to \emph{et vide si n=0}. The first English dual v is clear where the corresponding first French dual mark is absent; the other typed v's and raised 1's remain at their observed levels. English 74's many queries point individually to French 65: ``again'', ``by which of course we mean'', the G/Q gap, the sorites gloss, ``(only when)until'', the \emph{commode} gloss, and the two illegibility requests all remain. The points-of-$L$ phrase lacks the French raised o; the final $E$ index has a baseline 1 after lowered S. The capital C and ``in so far as'' replace a French hooked arrow and \emph{ssi}, respectively. The handwritten star/bullet indices on Grass and the coproduct sign before the indexed Grassmannians are absent in the English. Its unmatched N.B. bracket and final handwritten request to A.G. or Deligne are preserved. Sources: \href{https://grothendieck.umontpellier.fr/web/viewer.html?file=../24.pdf\#page=72}{cote 24, sites 71--74}; \href{https://grothendieck.umontpellier.fr/web/viewer.html?file=../25.pdf\#page=63}{cote 25, sites 62--65}.

\paragraph{The two final sheets (sites 75--76).}
The short handwritten French note on site 75 refers to having several important activities at once, one mathematical and another nonmathematical but worthwhile. Its initial \emph{Idem} and \emph{aussi} remain doubtful. The interlinear \emph{importantes} is connected to the gap after \emph{activités} by a large pen leader, whose right arm underlines the inserted word. Mirrored typing visible through the sheet belongs to the reverse and is described in a note rather than transcribed a second time. Site 76 is that separate typed CV ending, signed A. Grothendieck, with a statement about \emph{Survivre}, survival, and scientific activity, followed by a list of principal publications. Its two marginal pen brackets, separate short stroke, and folio number remain visible in the diplomatic transcription. The dark wedge below \emph{parmi} is a hole in the paper, established by the corresponding mirrored feature on site 75. This sheet's inclusion in the diplomatic volume preserves the archive sequence while its stray-sheet label explains its position. Sources: \href{https://grothendieck.umontpellier.fr/web/viewer.html?file=../24.pdf\#page=76}{cote 24, site 75}; \href{https://grothendieck.umontpellier.fr/web/viewer.html?file=../24.pdf\#page=77}{cote 24, site 76}.
